\section*{第 \ref{ch:b2:measurement-statistics} 章——化学测量的统计学}

\begin{solution}{exo:b2:measurement-statistics:1}
$\bar V = 62.41/5 = \qty{12.482}{mL}$; $\sum(V_i - \bar V)^2 = 0.00328$, $s = \sqrt{0.00328/4} \approx
\qty{0.029}{mL}$; $s/\sqrt5 \approx \qty{0.013}{mL}$.
\end{solution}

\begin{solution}{exo:b2:measurement-statistics:2}
$t = 2.776$（自由度为 4）：$12.482 \pm 2.776 \times 0.0128$，即 $(12.48 \pm 0.04)$~mL。
\end{solution}

\begin{solution}{exo:b2:measurement-statistics:3}
质量与体积的相对不确定度分别为 0.04~\% 和 0.06~\%；
按方和根合成，
\[\frac{u(c)}{c} = \sqrt{\Bigl(\frac{0.0002}{0.5000}\Bigr)^2 + \Bigl(\frac{0.15}{250.0}\Bigr)^2}
\approx 0.07~\%.\]
\end{solution}

\begin{solution}{exo:b2:measurement-statistics:4}
$\bar x = 1.5$, $\bar y = 0.31$, $S_{xx} = 5$, $S_{xy} = 0.99$: $b = 0.198$, $a = 0.31 - 0.198 \times 1.5
= 0.013$.
\end{solution}

\begin{solution}{exo:b2:measurement-statistics:5}
$|10.12 - 10.00|/(0.08/\sqrt5) = 0.12/0.036 \approx 3.4 > 2.776$：差异显著，
该方法有\omterm{def:b2:measurement-statistics:validation}{偏倚}（约为 \qty{+0.12}{mg/L}）。
\end{solution}

\begin{solution}{exo:b2:measurement-statistics:6}
$\mathrm{pH} = -\log_{10}c$, $\mathrm d\,\mathrm{pH}/\mathrm dc = -1/(c\ln10)$:
$u(\mathrm{pH}) = (u(c)/c)/\ln10 = 0.02/2.303 \approx 0.009$.
\end{solution}

\begin{solution}{exo:b2:measurement-statistics:7}
$x_{\mathrm{LOD}} = 3 \times 0.0012/0.0098 \approx \qty{0.37}{\micro g/L}$; $x_{\mathrm{LOQ}} = 10
\times 0.0012/0.0098 \approx \qty{1.2}{\micro g/L}$.
\end{solution}

\begin{solution}{exo:b2:measurement-statistics:8}
\omterm{def:b2:measurement-statistics:calibration}{标准曲线}给出测定溶液中的浓度为 \qty{3.11}{mg/L}；样品曾被稀释了五倍：
$3.11 \times 50.0/10.0 \approx \qty{15.6}{mg/L}$.
\end{solution}

\begin{solution}{exo:b2:measurement-statistics:9}
\qty{0.05}{mg/L} 为\omterm{def:b2:measurement-statistics:validation}{重复性}，\qty{0.15}{mg/L} 为\omterm{def:b2:measurement-statistics:validation}{再现性}。每个实验室都有
其自身的小\omterm{def:b2:measurement-statistics:validation}{偏倚}（校准、标准物质、仪器）；这些\omterm{def:b2:measurement-statistics:validation}{偏倚}因实验室而异，并增加
离散程度，因此\omterm{def:b2:measurement-statistics:validation}{再现性}总是更大。
\end{solution}

\begin{solution}{exo:b2:measurement-statistics:10}
$\partial S/\partial a = 0$ 即 $\sum(y_i - a - bx_i) = 0$：\omterm{def:b2:measurement-statistics:calibration}{残差}之和为零。两边同除以
$n$，得 $\bar y = a + b\bar x$，因此 $(\bar x, \bar y)$ 在直线上。
\end{solution}

\begin{solution}{exo:b2:measurement-statistics:11}
一步稀释：
\[\sqrt{(0.006/1.000)^2 + (0.08/100.0)^2} \approx 0.61~\%.\]
一次十倍稀释：
\[\sqrt{(0.02/10.00)^2 + (0.08/100.0)^2} \approx 0.22~\%,\]
两者按方和根合成得 $0.22\sqrt2 \approx 0.30~\%$。两步法的\omterm{def:b2:measurement-statistics:validation}{精密度}是一步法的两倍：小移液管主导了一步法
的不确定度。
\end{solution}

\begin{solution}{exo:b2:measurement-statistics:12}
$z = 1.6/\sqrt{0.4^2 + 0.5^2} \approx 2.5 > 2$：不相容；至少有一个实验室存在\omterm{def:b2:measurement-statistics:validation}{偏倚}。
取扩展不确定度（$k = 2$）时，第一个给出 $9.6 \pm 0.8$，包含 10；第二个给出
$11.2 \pm 1.0$，全部高于 10。对水样作任何判断之前，须先解决这一分歧，
例如借助标准参考物质。
% ledger: who:lead
\end{solution}

\begin{solution}{pb:b2:measurement-statistics:1}
\textbf{1.} 酸会改变原子化过程与响应；标准溶液与样品须具有相同的
基体，斜率方可适用。
\textbf{2.} $\bar x = 10$，$\bar y = 0.1004$，$S_{xx} = 250$，$S_{xy} = 2.445$：$b = 0.00978$ 每
\unit{\micro g/L}, $a = 0.1004 - 0.0978 = 0.0026$.
\textbf{3.} $+0.0004$, $-0.0005$, $+0.0006$, $-0.0013$, $+0.0008$.
\textbf{4.} 符号交替，无曲率：该直线在 0--\qty{20}{\micro g/L} 范围内适用。
\textbf{5.} $s_{y/x} = \sqrt{3.1 \times 10^{-6}/3} \approx 0.00102$.
\textbf{6.} 空白（酸、水、石墨炉）本身有较小的吸光度。
\textbf{7.} 灵敏度，即每 \unit{\micro g/L} 铅所对应的吸光度。
\textbf{8.} $\bar y_0 = 0.1010$; $s = 0.0030$.
\textbf{9.} $x_0 = (0.1010 - 0.0026)/0.00978 \approx \qty{10.06}{\micro g/L}$.
\textbf{10.} $s_{x_0} = (0.00102/0.00978)\sqrt{1/3 + 1/5 + (0.1010 - 0.1004)^2/(0.00978^2 \times 250)}
\approx 0.104 \times 0.730 \approx \qty{0.076}{\micro g/L}$.
\textbf{11.} $n - 2 = 3$; $t = 3.182$.
\textbf{12.} $10.06 \pm 3.182 \times 0.076 = (10.06 \pm 0.24)$~\unit{\micro g/L}.
\textbf{13.} $1/m$ 这一项由 1 降至 $1/3$：$s_{x_0}$ 约降低三分之一。
\textbf{14.} $f = 50.50/50.0 = 1.010$; $10.06 \times 1.010 \approx \qty{10.16}{\micro g/L}$.
\textbf{15.} $f = 1 + V_2/V_1$: $u^2(f) = (u(V_2)/V_1)^2 + (V_2u(V_1)/V_1^2)^2 = (0.005/50.0)^2 +
(0.50 \times 0.03/2500)^2$，$u(f) \approx 1.0 \times 10^{-4}$，相对值为 0.01~\%。
\textbf{16.} 否：0.01~\% 相对于 $0.076/10.06 \approx 0.75~\%$，可以忽略。
\textbf{17.} $s_{\text{空白}} \approx 0.00115$; $x_{\mathrm{LOD}} = 3 \times 0.00115/0.00978 \approx
\qty{0.35}{\micro g/L}$.
\textbf{18.} $10 \times 0.00115/0.00978 \approx \qty{1.2}{\micro g/L}$.
\textbf{19.} 是的，相差甚远。
\textbf{20.} 是：自来水中的测定值由 9.92 变至 \qty{10.40}{\micro g/L}。
\textbf{21.} $(0.104 - 0.0026)/0.00978 \times 1.010 \approx \qty{10.5}{\micro g/L}$，若单独看，
则会得出一个“高于指导值”的自信结论，而数据并不支持这一结论。
\textbf{22.} 构造区间所用的规则，在其所应用的
数据系列中，有 95~\% 包含真值；它并不是关于这一个值的概率。
\textbf{23.} 增加平行测定次数，并在 \qty{10}{\micro g/L} 附近增加标准溶液（即 $1/m$ 与
$1/n$ 两项），或采用灵敏度更高的方法；区间仅随 $1/\sqrt{\text{次数}}$ 缓慢缩小。
\textbf{24.} 通过分析有证参考水样，或在样品中加入已知量的
铅并检验回收率。
\textbf{25.} 在 \qty{10}{\micro g/L} 附近，常规方法的不确定度约为百分之几，如本例所示：
普通实验室无法可靠地检验更低的指导值。
\textbf{26.} $\boldsymbol{(10.16 \pm 0.24)}$~\unit{\micro g/L}（\omterm{def:b2:measurement-statistics:confidence}{置信水平} 95~\%）：区间包含
\qty{10}{\micro g/L}，因此本次测量既不能说明水样超过指导值，
也不能说明其符合指导值且留有余量。
\end{solution}
